% =====================================================================
%  §3 — A Semantic Model of Genetic Evidence
%
%  Core classes, dimensional vocabulary, conditional activation
%  mechanism. Written to be self-contained; subsequent sections
%  reference back to the table and the formal snippets here.
% =====================================================================

\section{A Semantic Model of Genetic Evidence}
\label{sec:model}

The model contains three parts: a small set of \emph{core classes}
(\cref{sec:model-classes}) that define what a unit of evidence is and
how evidence items compose; a \emph{dimensional vocabulary}
(\cref{sec:model-dimensions}) that describes each evidence item along
roughly twenty axes; and a \emph{conditional-activation} mechanism
(\cref{sec:model-conditional}) that keeps the vocabulary compact by
requiring dimensions only when they are semantically appropriate.

\subsection{Core Classes}
\label{sec:model-classes}

Three abstract classes form the core of the model:

\begin{description}
    \item[\ScientificEvidence] A unit of evidence, typically linked to a
          publication or dataset and described by a set of dimensions (see
          \cref{tab:dimensions-core,tab:dimensions-cond}). It carries an
          \texttt{assertion} property containing the object-level,
          source-anchored claims made by the source, for example that a
          variant is associated with a phenotype. These are the claims
          recorded in the pilot (\cref{sec:workflow}). A single paper may
          give rise to multiple \ScientificEvidence{} items when it reports
          distinct evidential claims, such as a case--control association and
          an \emph{in vivo} functional assay of the same variant.
    \item[\EvidenceVariable] A named quantity extracted from one or more
        \ScientificEvidence{} items, such as the gnomAD allele frequency of
        a variant, the LOD score in a family, or the odds ratio reported by
        a GWAS. Evidence variables are the terms from which predicates over
        Evidence are constructed. They are also described by dimensions,
        because the same quantity, such as an odds ratio, may carry different
        epistemic weight depending on how it was obtained.
    \item[\EvidenceAssertion] A meta-predicate, that is, a predicate about a
        predicate. It determines whether \EvidenceVariable{}s are used
        appropriately in a predicate over Evidence, namely in a decision
        rule. This class is the model's point of contact with the reasoning
        layer described in~\cite{our_preprint}, where such meta-predicates
        assert which evidence a rule may use. This paper constructs no
        decision rules and therefore instantiates no \EvidenceAssertion{}.
        The class remains part of the model, but the interpretability and
        AI-readiness of the reasoning layer are the primary focus
        of~\cite{our_preprint}, not of the present pilot.
\end{description}

The \texttt{assertion} property and the \EvidenceAssertion{} class share
a name but denote different things. The property records the object-level
claims captured in the pilot, whereas the class is a meta-predicate over
any decision rule that would use those claims. The property name is
retained for compatibility with the released annotations and will be
renamed, for example to \texttt{selector}, in a future schema revision.

Specialising these classes for the genetic domain yields
\GeneticEvidence, \GeneticEvidenceVariable, and
\GeneticEvidenceAssertion. These specialisations inherit the provenance
requirement and add the dimensional annotations described next.

\subsection{Dimensions}
\label{sec:model-dimensions}

Each \GeneticEvidence{} instance is described using a dimensional
vocabulary summarised in \cref{tab:dimensions-core}, which lists the
always-required dimensions, and \cref{tab:dimensions-cond}, which gives
representative conditional dimensions.\footnote{The full enumeration of
conditional dimensions, including those omitted from
\cref{tab:dimensions-cond} for space, is maintained in the project
repository at
\url{https://github.com/ForomePlatform/genetic-evidence-model/blob/master/schema/dimensions.md}.}

Six dimensions are \emph{always required}: Knowledge Domain, Method,
Target Type, Resolution, Credibility, and Phenotype Scale. No
well-formed genetic-evidence item is fully specified without them.
Together, they tell a downstream reader what kind of evidence the item
is, how it was obtained, what it concerns, and how credible it is in
context.

The dimensions differ in internal structure, and each is modeled
according to its own semantics rather than forced into a single
hierarchy. Knowledge Domain values are sibling framings rather than a
subsumption hierarchy. A single item may carry several values at once,
as when a functional study is interpreted in a clinical-genetics
context; the multi-valued cardinality captures this directly.

Phenotype Scale is an ordinal scale, from molecular to clinical, rather
than a class hierarchy. Modeling it through subsumption would therefore
misrepresent the dimension. Method, by contrast, is genuinely
hierarchical, and its \texttt{is\_a} relations are encoded in the SHACL
distribution. The chosen structure is therefore part of the definition
of each dimension.

Phenotype Scale and Variant Ascertainment were added during the annotation
work (\cref{sec:workflow}). Phenotype Scale distinguishes molecular
phenotypes, such as binding or activity, from cellular, histological,
organismal, and clinical phenotypes. It is independent of the scale of the
genetic target.

Variant Ascertainment records how a variant entered a study. It is meaningful
only for variant-level targets and is therefore a \emph{conditional}
dimension, required when Target Type is \dimval{Variant}. Both dimensions
are revisited in \cref{sec:discussion}.

Credibility records the methodological soundness of an asserted finding, and
nothing else. It is a structured decomposition rather than an open key-value
field: one overall ordinal rating together with separately rated facets for
statistical power or sample size, independent replication, multiple-testing
control, ancestry or population-stratification control, and ascertainment.
The facets remain separate rather than being collapsed into a single score
because clinical and research users may weigh them differently.

This separation is inspired by GRADE's general practice of assessing the
factors that affect certainty separately~\cite{grade2008}. GEM does not
adopt the GRADE framework or its rating rules, which were developed for
bodies of clinical evidence, particularly randomized and observational
studies. The GEM facets instead support credibility judgments about
individual evidence items from basic and pre-clinical research.

The overall rating is anchored on SEPIO confidence, defined as the degree of
belief that an asserted proposition is true~\cite{sepio}, and is implemented
as a genetics-specific value set bound to that attribute, as prescribed by
SEPIO profiles (Supplementary Table~\ref*{supp:tab:credibility-facets}).

Two constructs are excluded from credibility by design. First, the
direction or hedging of a finding belongs to the assertion itself and is
represented by the width of the value range it constrains. An honestly
reported null result from a rigorous study may therefore have high
credibility. Second, the extent to which evidence bears on a particular
question is applicability. It is captured by Relevance and Specificity and
computed for each question. Combining methodological soundness, assertion
content, and applicability in one field would conflate distinct constructs.

The remaining dimensions are \emph{conditional}: they are required, or
even meaningful, only when other dimensions take particular values. For
example, mode of inheritance is not meaningful for a polygenic score, and
knockout type is not meaningful for a human case--control study. Rather
than marking such dimensions as ``optional'' and relying on convention, the
model assigns explicit activation conditions to each, as described in
\cref{sec:model-conditional}. These conditions specify when a dimension
must be present; otherwise, it may be omitted. This makes the requirement
status transparent to both human curators and automated validators.

\begin{table}[tbp]
\centering
\caption{Always-required dimensions of the \GeneticEvidence{} model.
These apply to every evidence item regardless of context. Phenotype Scale
was added to this set during the annotation work
(\cref{sec:workflow}). KD abbreviates Knowledge Domain.}
\label{tab:dimensions-core}
\small
\begin{tabularx}{\linewidth}{@{}llllX@{}}
\toprule
Dimension              & Value type        & Cardinality & Objectivity & Example values \\
\midrule
Knowledge Domain       & categorical       & multiple & obj.       & \dimval{Human Genetics}, \dimval{Population Genetics}, \dimval{Model Organism}, \dimval{Gene Function} \\
Method                 & categorical       & multiple & obj.       & \dimval{Statistical Genetics}, \dimval{In Vivo}, \dimval{In Vitro}, \dimval{Clinical Evidence} \\
Target Type            & categorical       & single   & obj.       & \dimval{Gene}, \dimval{Variant}, \dimval{Interval}, \dimval{Transcript} \\
Resolution             & categorical       & single   & obj.       & \dimval{Window}, \dimval{Gene}, \dimval{Functional Element}, \dimval{Position}, \dimval{Variant} \\
Credibility            & rating ${+}$ facets & multiple & subj.      & overall rating ${+}$ \texttt{sample\_size}, \texttt{replication}, \texttt{multiple\_testing}, \texttt{stratification\_control}, \ldots \\
Phenotype Scale        & categorical       & single   & obj.       & \dimval{Molecular}, \dimval{Cellular}, \dimval{Histological}, \dimval{Organismal}, \dimval{Clinical} \\
\bottomrule
\end{tabularx}
\end{table}

\begin{table}[tbp]
\centering
\caption{Conditional dimensions of the \GeneticEvidence{} model.
Each is required only when its activation condition holds. KD, MT, and TT
abbreviate Knowledge Domain, Method, and Target Type, respectively;
$\supseteq$ denotes set membership for multi-valued dimensions.}
\label{tab:dimensions-cond}
\small
\begin{tabularx}{\linewidth}{@{}lllX@{}}
\toprule
Dimension              & Value type   & Cardinality & Activation condition \\
\midrule
Mode of Inheritance    & categorical  & single   & KD $\supseteq$ \dimval{Human Genetics} AND TT = \dimval{Gene} \\
Mendelian Segregation  & boolean      & single   & KD $\supseteq$ \dimval{Human Genetics} AND TT = \dimval{Gene} \\
Penetrance             & categorical  & single   & KD $\supseteq$ \dimval{Human Genetics} AND TT = \dimval{Variant} \\
Variant Ascertainment  & categorical  & multiple & TT = \dimval{Variant} \\
Measurement Target     & categorical  & multiple & KD $\supseteq$ \dimval{Gene Function} \\
Gene Relation          & categorical  & single   & KD $\supseteq$ \dimval{Gene Function} \\
Organism               & categorical  & multiple & MT $\supseteq$ \dimval{In Vivo} OR KD $\supseteq$ \dimval{Model Organism} \\
Knockout Type          & categorical  & single   & KD $\supseteq$ \dimval{Model Organism} \\
\bottomrule
\end{tabularx}
\end{table}

\subsection{Conditional Activation}
\label{sec:model-conditional}

Conditional activation is the least conventional part of the model, but it
keeps the vocabulary compact without sacrificing expressive power. Rather
than using a fixed schema in which every \GeneticEvidence{} item has every
dimension, with many dimensions marked as optional and left empty, the model
requires particular dimensions only when specified conditions over other
dimensions hold. This mirrors expert reasoning. A curator reading a paper
about an X-linked recessive pedigree asks about mode of inheritance because
the evidence makes that question meaningful, not because a global schema
requires it.

The activation conditions are encoded as SHACL shapes with SPARQL-based
constraints. This is the standard SHACL mechanism for expressing
requirements that cannot be captured by simple cardinality constraints on a
single property. Two representative shapes illustrate the pattern. For
brevity, we show only the \texttt{sh:sparql} body of each and omit the
outer NodeShape wrapper. The full shapes, together with the base
\texttt{GeneticEvidenceShape} and the simpler
\texttt{variant\_ascertainment} rule, are available in the project
repository with an annotated walkthrough.\footnote{\url{https://github.com/ForomePlatform/genetic-evidence-model/blob/master/schema/examples.md}}

Checking activation conditions is a closed-world question: whether an
annotation contains every dimension required by the evidence item. SHACL can
answer this question directly. OWL's open-world semantics cannot express
this form of completeness, which is one reason the constraints are encoded
in SHACL rather than as ontology axioms.

The first example encodes a conjunction across two classification axes.
\emph{Mode of Inheritance} is required when Knowledge Domain contains
\dimval{Human Genetics} \textbf{and} Target Type is \dimval{Gene}
(\cref{lst:shacl-moi}).

\begin{lstlisting}[language=yaml, caption={Compact form of the conditional-activation shape for \emph{Mode of Inheritance}. The outer NodeShape wrapper is omitted. The two triple patterns in \texttt{WHERE} form the conjunction; \texttt{FILTER NOT EXISTS} identifies instances that satisfy the condition but lack the required dimension.}, label={lst:shacl-moi}]
sh:sparql [ sh:select """
  SELECT $this WHERE {
    $this gem:knowledgeDomain gem:HUMAN_GENETICS .
    $this gem:targetType      gem:GENE .
    FILTER NOT EXISTS { $this gem:modeOfInheritance ?m }
  }""" ] .
\end{lstlisting}

The second example encodes a disjunction across two classification axes.
\emph{Organism} is required when Method contains \dimval{In Vivo}
\textbf{or} Knowledge Domain contains \dimval{Model Organism}
(\cref{lst:shacl-org}). The two disjuncts refer to different dimensions:
one is a Method value and the other is a Knowledge Domain value. This
illustrates that activation conditions can span the dimensional vocabulary
rather than being confined to a single axis.

\begin{lstlisting}[language=yaml, caption={Compact form of the conditional-activation shape for \emph{Organism}. The disjunction is expressed by a SPARQL \texttt{UNION} whose two branches refer to different dimensions.}, label={lst:shacl-org}]
sh:sparql [ sh:select """
  SELECT $this WHERE {
    { $this gem:method          gem:IN_VIVO . }
    UNION
    { $this gem:knowledgeDomain gem:MODEL_ORGANISM . }
    FILTER NOT EXISTS { $this gem:organism ?o }
  }""" ] .
\end{lstlisting}

In each case, the \texttt{SELECT} returns every \GeneticEvidence{}
instance that satisfies the activation condition but lacks the required
dimension. The validator reports these instances as constraint violations
with a localized diagnostic. Writing each conditional rule explicitly adds
some complexity, but it improves downstream clarity: a reader of an
annotation file can see why a particular dimension is required or absent,
and a validator can check conformance automatically.

\Cref{sec:workflow} returns to conditional activation in the context of
the annotation workflow.
